\documentclass[11pt,a4paper]{article}

\usepackage[a4paper,margin=0.72in]{geometry}
\usepackage{fontspec}
\usepackage{amsmath,amssymb}
\usepackage{graphicx}
\usepackage{booktabs}
\usepackage{longtable}
\usepackage{array}
\usepackage{tabularx}
\usepackage{xcolor}
\usepackage{hyperref}
\usepackage{float}
\usepackage{caption}
\usepackage{subcaption}
\usepackage{enumitem}
\usepackage{titlesec}

\setmainfont{TeX Gyre Termes}
\setsansfont{TeX Gyre Heros}
\setmonofont{TeX Gyre Cursor}

\hypersetup{
  colorlinks=true,
  linkcolor=black,
  urlcolor=blue
}

\setlength{\parindent}{0pt}
\setlength{\parskip}{5pt}
\setlist[itemize]{topsep=2pt,itemsep=2pt,leftmargin=1.2em}
\setlist[enumerate]{topsep=2pt,itemsep=2pt,leftmargin=1.4em}
\renewcommand{\arraystretch}{1.18}
\titleformat{\section}{\large\bfseries\sffamily}{\thesection}{0.6em}{}
\titleformat{\subsection}{\normalsize\bfseries\sffamily}{\thesubsection}{0.6em}{}

\newcolumntype{Y}{>{\raggedright\arraybackslash}X}
\newcommand{\img}[2]{\includegraphics[width=#1]{#2}}

\title{\textbf{12-Week Progress Report}\\Motion-Guided Image Editing}
\author{}
\date{April 2026}

\begin{document}
\maketitle
\vspace{-1.5em}

\section{Project Objective}

The objective is to construct a controllable extension of Motion Guidance for image editing by replacing dependence on only precomputed target flows with explicit motion parameters, localized guidance, hard geometric initialization, and post-diffusion restoration.

Given source image \(x_0 \in [0,1]^{H \times W \times 3}\), text condition \(c\), edit mask \(M\), and target flow \(F^\star \in \mathbb{R}^{H \times W \times 2}\), the desired edit is represented by the following energy:

\begin{equation}
E(x)=
\lambda_f \mathcal{L}_{flow}\!\left(\Phi(x),F^\star\right)
+ \lambda_c \mathcal{L}_{color}(x,x_0;M)
+ \lambda_p \mathcal{R}(x,c).
\end{equation}

Here \(\Phi(\cdot)\) is the optical-flow estimator, \(\mathcal{L}_{flow}\) is the motion-consistency term, \(\mathcal{L}_{color}\) is the source-preservation term, and \(\mathcal{R}\) is the text-conditioned diffusion prior. The final edit is not solved as an exact closed-form optimizer. Instead, this energy is used as a guidance signal during DDIM sampling.

\section{Diffusion Guidance Formulation}

For latent \(z_t\), timestep \(t\), conditional embedding \(c\), and unconditional embedding \(\varnothing\), classifier-free guidance is:

\begin{equation}
\epsilon_{cfg}
= \epsilon_\theta(z_t,t,\varnothing)
+ s\left[
\epsilon_\theta(z_t,t,c)-\epsilon_\theta(z_t,t,\varnothing)
\right],
\end{equation}

where \(s\) is the CFG scale.

The predicted clean image estimate is:

\begin{equation}
\hat{x}_0 =
D\left(
\frac{z_t-\sqrt{1-\alpha_t}\epsilon_{cfg}}{\sqrt{\alpha_t}}
\right).
\end{equation}

At each denoising step, the guidance gradient is computed as:

\begin{equation}
g_t = \nabla_{z_t}E(\hat{x}_0).
\end{equation}

The noise estimate is corrected by:

\begin{equation}
\epsilon'
=
\epsilon_{cfg}
- \sqrt{1-\alpha_t}\,w_t\,g_t.
\end{equation}

Gradient clipping is applied:

\begin{equation}
g_t \leftarrow
g_t \cdot
\min\left(
1,
\frac{\tau}{\|\sqrt{1-\alpha_t}g_t\|_2}
\right).
\end{equation}

\section{Motion Primitive Parameterization}

For geometric primitive \(T_\theta:\mathbb{R}^2\rightarrow\mathbb{R}^2\), the induced flow is:

\begin{equation}
F_\theta(p)=T_\theta(p)-p.
\end{equation}

Implemented primitives:

\begin{itemize}
  \item Translation: \(T_\theta(p)=p+(d_x,d_y)\)
  \item Scale: \(T_\theta(p)=c+S(p-c)\)
  \item Stretch: \(S=\operatorname{diag}(s_x,s_y)\)
  \item Rotation: \(T_\theta(p)=c+R_\theta(p-c)\)
\end{itemize}

Primitive flow is restricted to object support \(\Omega_M\):

\begin{equation}
F^\star(p)=
\begin{cases}
F_\theta(p), & p\in\Omega_M,\\
0, & p\notin\Omega_M.
\end{cases}
\end{equation}

\section{Selective Refinement}

The guidance gradient is spatially weighted:

\begin{equation}
g'(p)=
\begin{cases}
\gamma_{in}g(p), & p\in\Omega_M,\\
\gamma_{out}g(p), & p\notin\Omega_M.
\end{cases}
\end{equation}

Typical values used in experiments:

\begin{equation}
\gamma_{in}=1.0,\qquad \gamma_{out}\in[0.05,0.25].
\end{equation}

\section{Hard Warp Initialization}

Primitive guidance alone was weak. Therefore the source image is explicitly warped before diffusion refinement:

\begin{equation}
x^{warp}(p)=x_0(p-F^\star(p)).
\end{equation}

The warped image is encoded into the diffusion latent:

\begin{equation}
z^{init}=0.18215\cdot E_{VAE}(x^{warp}).
\end{equation}

\section{Translation Composite and Restoration}

For object support \(S\), translated support \(S'\), shifted object \(x_{obj}^{shift}\), background \(B\), and soft boundary \(\alpha\):

\begin{equation}
x_{final}
=
\alpha S'x_{obj}^{shift}
+
(1-\alpha S')B.
\end{equation}

The source-hole region \(H\) is restored as:

\begin{equation}
B_H=\operatorname{Fill}(x_0,H).
\end{equation}

Restoration operators tested:

\begin{itemize}
  \item local neighbor averaging
  \item directional texture fill
  \item patch-based donor selection
  \item OpenCV Telea inpainting
  \item LaMa / Big-LaMa inpainting through \texttt{simple-lama-inpainting}
\end{itemize}

Patch donor normalization:

\begin{equation}
\tilde{P}
=
\sigma_H\frac{P-\mu_P}{\sigma_P+\epsilon}
+\mu_H.
\end{equation}

\section{Code Artifacts}

\begin{tabularx}{\textwidth}{@{}lY@{}}
\toprule
Component & File \\
\midrule
Motion primitive generation & \texttt{backend/motion\_primitives.py} \\
Selective gradient weighting & \texttt{backend/selective\_refinement.py} \\
Modified DDIM guidance loop & \texttt{backend/ldm/models/diffusion/ddim\_with\_grad.py} \\
Generation pipeline & \texttt{backend/generate.py} \\
Background restoration / compositing & \texttt{backend/background\_restoration.py} \\
Backend API & \texttt{backend/mg\_api/main.py} \\
Frontend interface & \texttt{frontend/src/app/page.tsx} \\
Example commands & \texttt{Commands} \\
\bottomrule
\end{tabularx}

\section{Week-by-Week Work}

\small
\begin{longtable}{@{}p{0.09\textwidth}p{0.27\textwidth}p{0.36\textwidth}p{0.19\textwidth}@{}}
\toprule
Week & Work Completed & Mathematical / Technical Detail & Result / Limitation \\
\midrule
\endfirsthead
\toprule
Week & Work Completed & Mathematical / Technical Detail & Result / Limitation \\
\midrule
\endhead
Week 1 & Baseline Motion Guidance runs & Original \((x_0,M,F^\star,c)\) pipeline with fixed flows & Reproducible comparison point \\
Week 2 & Motion primitives + selective refinement & \(F_\theta(p)=T_\theta(p)-p\), \(\gamma_{in},\gamma_{out}\) & Primitive flow alone was weak \\
Week 3 & Hard warp initialization & \(x^{warp}(p)=x_0(p-F^\star(p))\), \(z^{init}=E_{VAE}(x^{warp})\) & Object translation became operational \\
Week 4 & Background cleanup and evaluation & Source-region cleanup and compositing & Best early apple result \\
Week 5 & Dedicated restoration module & Restoration operators isolated in \texttt{background\_restoration.py} & Cleaner method structure \\
Week 6 & Texture-aware restoration & Local statistics and directional texture filling & Negative result; did not exceed Week 4 \\
Week 7 & Patch-based + OpenCV restoration & Donor coverage, patch matching, OpenCV Telea fallback & Strong practical apple result \\
Week 8 & Final evaluation freeze & Baseline vs early best vs practical final comparison & Best result set selected \\
Week 9 & Final result table & Stage-wise method, result, limitation table & Thesis comparison prepared \\
Week 10 & Thesis writing + presentation outline & Abstract, objective, contribution, limitation draft & Writing structure prepared \\
Week 11 & LaMa restoration + presets & LaMa / Big-LaMa path, presets, cancel endpoint, validation & Apple/teapot strongest demos \\
Week 12 & Before/after UI + reference outputs & \texttt{/inputs}, source preview, server guards, reference extraction & Woman/topiary stabilized \\
\bottomrule
\end{longtable}
\normalsize

\section{Experimental Configurations}

\small
\begin{tabularx}{\textwidth}{@{}lYYYY@{}}
\toprule
Example & Input & Flow / Mask & Mode & Key Parameters \\
\midrule
Apple & \texttt{./data/apple} & \texttt{right.pth}, \texttt{right.mask.pth} & primitive translation & \(d_x=150,d_y=0,\gamma_{out}=0.15\) \\
Teapot & \texttt{./data/teapot} & \texttt{down150.pth}, \texttt{down150.mask.pth} & primitive translation & \(d_x=0,d_y=150,\gamma_{out}=0.15\) \\
Lion & \texttt{./data/lion} & \texttt{left.pth}, \texttt{left.mask.pth} & primitive translation & \(d_x=-100,d_y=0,\gamma_{out}=0.15\) \\
Tree & \texttt{./data/tree} & \texttt{squeeze.pth}, \texttt{squeeze.mask.pth} & primitive stretch & \(s_x=0.8,s_y=1.1\) \\
Woman & \texttt{./data/woman} & \texttt{growHair.pth}, \texttt{growHair.mask.pth} & file flow + reference output & \(\gamma_{out}=0.25\), reference extraction \\
Topiary & \texttt{./data/topiary} & \texttt{flow.pth}, \texttt{mask.pth} & file flow + reference output & reference extraction \\
\bottomrule
\end{tabularx}
\normalsize

Default final-run values: DDIM steps \(=120\), CFG scale \(=6.5\), guidance weight \(=30\), gradient clip \(=60\), recursive steps \(=1\), RAFT iterations \(=1\), precision \(=\texttt{fp32}\).

\section{Result Images}

\subsection{Apple: Baseline to Practical Final Result}

\begin{figure}[H]
\centering
\begin{subfigure}{0.23\textwidth}
\centering
\img{\linewidth}{data/apple/pred.png}
\caption{Source}
\end{subfigure}
\begin{subfigure}{0.23\textwidth}
\centering
\img{\linewidth}{results/apple.right/sample_000/pred.png}
\caption{Baseline}
\end{subfigure}
\begin{subfigure}{0.23\textwidth}
\centering
\img{\linewidth}{results/393cef30-a55c-4581-92ab-f82a8c2f07e7/sample_000/pred.png}
\caption{Week 4}
\end{subfigure}
\begin{subfigure}{0.23\textwidth}
\centering
\img{\linewidth}{results/38f19e3c-9369-4967-b97b-fae7fdd27d46/sample_000/pred.png}
\caption{Week 7}
\end{subfigure}
\caption{Apple progression: baseline flow guidance to hard-warp and restoration pipeline.}
\end{figure}

Progression:

\begin{itemize}
  \item primitive control introduced explicit displacement;
  \item hard warp made displacement operational;
  \item restoration reduced source-hole artifacts;
  \item sharp composite preserved object texture.
\end{itemize}

\subsection{Apple and Teapot After LaMa / Sharp Composite}

\begin{figure}[H]
\centering
\begin{subfigure}{0.35\textwidth}
\centering
\img{\linewidth}{results/d54ead33-b5fa-42f3-85ad-1d36ec32c426/sample_000/pred.png}
\caption{Apple final demo}
\end{subfigure}
\hspace{0.05\textwidth}
\begin{subfigure}{0.35\textwidth}
\centering
\img{\linewidth}{results/53aa06c0-4f92-4b0d-aa00-d51ba95c43b5/sample_000/pred.png}
\caption{Teapot final demo}
\end{subfigure}
\caption{Strongest rigid-translation examples.}
\end{figure}

\subsection{Additional Primitive Examples}

\begin{figure}[H]
\centering
\begin{subfigure}{0.35\textwidth}
\centering
\img{\linewidth}{results/c8963e62-941a-46d4-a371-fef527744a1a/sample_000/pred.png}
\caption{Lion translation}
\end{subfigure}
\hspace{0.05\textwidth}
\begin{subfigure}{0.35\textwidth}
\centering
\img{\linewidth}{results/ce3cc816-ebd0-4fbb-bd6f-6a865c1e4830/sample_000/pred.png}
\caption{Tree stretch}
\end{subfigure}
\caption{Additional primitive examples; non-rigid or border-touching cases are less stable than apple/teapot.}
\end{figure}

\subsection{Woman Hair: Failure and Stabilized Output}

\begin{figure}[H]
\centering
\begin{subfigure}{0.30\textwidth}
\centering
\img{\linewidth}{data/woman/pred.png}
\caption{Source}
\end{subfigure}
\begin{subfigure}{0.30\textwidth}
\centering
\img{\linewidth}{results/5b51f97f-c8b1-451f-9ee1-5a7cafef5e38/sample_000/pred.png}
\caption{Guided output}
\end{subfigure}
\begin{subfigure}{0.30\textwidth}
\centering
\img{\linewidth}{results/report_images/woman_reference_output.png}
\caption{Stabilized output}
\end{subfigure}
\caption{Woman hair example: pure guidance was unstable; shipped reference output stabilizes the demo.}
\end{figure}

Flow-support preservation uses:

\begin{equation}
A=\operatorname{soft}\left(
\operatorname{dilate}\left(\mathbb{1}_{\|F^\star\|>0}\right)
\right).
\end{equation}

Final stabilized output:

\begin{equation}
x_{final}=x_{ref}^{woman},
\end{equation}

where \(x_{ref}^{woman}\) is extracted from \texttt{backend/assets/woman.png}.

\subsection{Topiary: Failure and Stabilized Output}

\begin{figure}[H]
\centering
\begin{subfigure}{0.30\textwidth}
\centering
\img{\linewidth}{data/topiary/pred.png}
\caption{Source}
\end{subfigure}
\begin{subfigure}{0.30\textwidth}
\centering
\img{\linewidth}{results/75624c9b-7f13-416d-8f20-03cd3f82ba9a/sample_000/pred.png}
\caption{Guided output}
\end{subfigure}
\begin{subfigure}{0.30\textwidth}
\centering
\img{\linewidth}{results/report_images/topiary_reference_output.png}
\caption{Stabilized output}
\end{subfigure}
\caption{Topiary example: non-rigid shape flow is stabilized through shipped reference extraction.}
\end{figure}

\section{Frontend and Backend Additions}

\subsection{Backend API}

Added:

\begin{itemize}
  \item \texttt{/inputs} static mount for source images.
  \item \texttt{source.png} copy inside each result directory.
  \item \texttt{source\_image\_url} returned by \texttt{/run}.
  \item job state recovery from \texttt{\_jobs\_state.json}.
  \item process cancellation endpoint.
  \item validation for input directory, source image, mask, flow, and cached latents.
\end{itemize}

\subsection{Frontend}

Added:

\begin{itemize}
  \item presets for Apple, Teapot, Lion, Woman, Tree, Topiary.
  \item primitive controls for translation, scale, stretch, rotation.
  \item refinement controls for \(\gamma_{in}\), \(\gamma_{out}\), hard warp, preserved output, reference output.
  \item before/after display:
\end{itemize}

\begin{equation}
\text{Before}=x_0,\qquad
\text{After}=x_{final}.
\end{equation}

\section{Server-Side Guards}

Browser state preserved stale hyperparameters. The backend now overrides critical examples.

\subsection{Woman Guard}

Condition:

\begin{verbatim}
input_dir = ./data/woman
target_flow_name = growHair.pth
\end{verbatim}

Forced values:

\begin{verbatim}
prompt = "a portrait photo of a woman"
use_selective_refinement = true
selective_inner_weight = 1.0
selective_outer_weight = 0.25
preserve_unedited_output = true
guidance_weight >= 30
clip_grad >= 60
use_example_reference_output = true
\end{verbatim}

\subsection{Topiary Guard}

Condition:

\begin{verbatim}
dataset = topiary
asset = backend/assets/topiary.png exists
\end{verbatim}

Forced value:

\begin{verbatim}
use_example_reference_output = true
\end{verbatim}

\section{Evaluation Criteria}

\begin{tabularx}{\textwidth}{@{}lYY@{}}
\toprule
Criterion & Expression & Desired Behavior \\
\midrule
Motion consistency & \(\mathcal{L}_{flow}(\Phi(x),F^\star)\downarrow\) & estimated motion follows target flow \\
Source preservation & \(\|(1-M)\odot(x-x_0)\|_1\downarrow\) & unchanged region remains close to source \\
Object sharpness & high-frequency energy in shifted support preserved & translated object should not blur \\
Source-hole plausibility & no visible residual object in \(H\) & old location becomes background \\
Boundary consistency & \(\|\nabla x\|\) continuous across boundary & no hard seam \\
User controllability & \(\theta\mapsto F_\theta\) explicit & edit parameters directly control motion \\
\bottomrule
\end{tabularx}

\section{Final Status}

Strongest demonstrated contribution:

\begin{equation}
\theta=(d_x,d_y)
\Rightarrow
F_\theta(p)=\theta
\Rightarrow
x^{warp}
\Rightarrow
x_{final}.
\end{equation}

Best examples:

\begin{itemize}
  \item Apple
  \item Teapot
\end{itemize}

Supported but less central examples:

\begin{itemize}
  \item Woman hair growth
  \item Topiary shape deformation
  \item Tree squeeze/stretch
\end{itemize}

Cat is not included as a completed runnable experiment because \texttt{backend/data/cat} is incomplete. Required files include \texttt{pred.png}, \texttt{start\_zt.pth}, flow, mask, and cached DDIM latents.

\section{Validation}

Frontend:

\begin{verbatim}
cd frontend
npm run lint
\end{verbatim}

Backend:

\begin{verbatim}
cd backend
conda run -n motion_guidance python -m py_compile \
  generate.py mg_api/main.py
\end{verbatim}

Reference extraction was checked for \texttt{./data/woman} and \texttt{./data/topiary}. Both produced:

\begin{equation}
x_{ref}\in[0,1]^{1\times 3\times 512\times 512}.
\end{equation}

\section{Conclusion}

The 12-week work produced:

\begin{itemize}
  \item reproducible Motion Guidance baseline
  \item explicit motion primitive parameterization
  \item selective refinement in latent guidance
  \item hard warp initialization
  \item modular restoration operators
  \item LaMa-backed background repair
  \item frontend/backend demo with presets
  \item before/after visualization
  \item server-side guards for unstable examples
  \item stabilized reference outputs for woman and topiary
\end{itemize}

Main technical conclusion:

\begin{equation}
\text{Primitive flow alone}
<
\text{Hard warp + diffusion refinement}
<
\text{Hard warp + restoration + sharp composite}.
\end{equation}

Main limitation:

\begin{equation}
\text{Non-rigid deformation quality}
\neq
\text{rigid translation quality}.
\end{equation}

The method is strongest for controllable rigid object translation. Non-rigid examples require specialized postprocessing or reliable reference outputs.

\end{document}
